\documentclass[11pt]{article}

\usepackage[T1]{fontenc}
\usepackage{amsmath,amssymb,amsthm}
\usepackage{geometry}
\usepackage[colorlinks=true,linkcolor=blue,citecolor=blue,urlcolor=blue]{hyperref}
\geometry{margin=1in}

\newtheorem{theorem}{Theorem}
\newtheorem{corollary}{Corollary}
\newtheorem{definition}{Definition}
\newtheorem{remark}{Remark}

\newcommand{\dist}{\operatorname{dist}}
\newcommand{\WAP}{\mathrm{WAP}}
\newcommand{\eps}{\varepsilon}

\title{A Bochner \(L^p\) lifting lemma for weakly compact approximation}
\author{Run fa\_banach\_001, agent lane 4}
\date{2026-06-28}

\begin{document}
\maketitle

\begin{abstract}
Odell and Tylli ask whether \(L^p([0,1];E)\) has the weakly compact approximation property whenever \(E\) has it and \(1<p<\infty\).  We prove a partial permanence theorem: the answer is affirmative on every weakly compact subset of \(L^p(\mu,E)\) which is \(L^p\)-tight with respect to weakly compact subsets of \(E\).  The proof is a direct lifting argument.  A weakly compact operator \(S:E\to E\) approximating the identity on the controlling weakly compact set \(W\subset E\) induces a weakly compact Bochner operator \(\widetilde S\) on \(L^p(\mu,E)\), and the \(L^p\)-tightness bounds the error.  This gives a clean reduction of the W.A.P. half of Problem 5.7 to the missing Bochner tightness step, and it also covers uniformly bounded \(\delta\mathcal S\)-sets.
\end{abstract}

\section{Source question}

Odell and Tylli define the weakly compact approximation property (W.A.P.) as follows.  A Banach space \(E\) has the W.A.P. if there is \(C<\infty\) such that for every weakly compact set \(D\subset E\) and every \(\eps>0\), there is a weakly compact operator \(V:E\to E\) satisfying
\[
  \sup_{x\in D}\|x-Vx\|<\eps,\qquad \|V\|\le C.
\]
In Problem 5.7 they ask:
\[
\textit{Does }L^p([0,1];E)\textit{ have the W.A.P. whenever }E\textit{ has the W.A.P. and }1<p<\infty?
\]
They also ask the analogous question for the inner W.A.P.  The present note addresses only the W.A.P. side, and only under an explicit tightness hypothesis on the weakly compact set being approximated.

\section{A tightness condition}

\begin{definition}
Let \((\Omega,\Sigma,\mu)\) be a probability space, \(1<p<\infty\), and \(E\) a Banach space.  A bounded set \(D\subset L^p(\mu,E)\) is called \(L^p\)-weakly-compactly tight if for every \(\eta>0\) there is a weakly compact set \(W\subset E\) such that
\[
  \sup_{f\in D}\left(\int_\Omega \dist(f(\omega),W)^p\,d\mu(\omega)\right)^{1/p}<\eta .
\]
\end{definition}

This condition is stronger than relative weak compactness in general.  It is exactly the extra range-control needed by the simple operator-lifting argument below.

\section{Proof intuition}

The natural way to lift a weakly compact approximant \(S:E\to E\) is pointwise:
\[
  (\widetilde S f)(\omega)=S(f(\omega)).
\]
The obstacle is that \(S\) only approximates the identity on a weakly compact subset of \(E\), while a Bochner function may take values outside that subset.  The \(L^p\)-tightness hypothesis says that, for the particular weakly compact family \(D\), those values are close in \(L^p\)-mean to a weakly compact \(W\subset E\).  Once \(S\) approximates the identity on \(W\), the pointwise estimate
\[
  \|x-Sx\|\le \sup_{w\in W}\|w-Sw\|+(1+\|S\|)\dist(x,W)
\]
integrates to the required uniform \(L^p\)-estimate.

The lifted operator is weakly compact because every weakly compact operator factors through a reflexive space.  If \(S=BA\) with \(A:E\to R\), \(B:R\to E\), and \(R\) reflexive, then
\[
  \widetilde S = \widetilde B\,\widetilde A
\]
factors through \(L^p(\mu,R)\), which is reflexive for \(1<p<\infty\).

\section{The lifting theorem}

\begin{theorem}\label{thm:tight-lift}
Let \(1<p<\infty\), let \((\Omega,\Sigma,\mu)\) be a probability space, and let \(E\) have the W.A.P. with constant \(C\).  Let \(D\subset L^p(\mu,E)\) be a weakly compact \(L^p\)-weakly-compactly tight set.  Then for every \(\eps>0\) there is a weakly compact operator \(T:L^p(\mu,E)\to L^p(\mu,E)\) such that
\[
  \|T\|\le C,\qquad \sup_{f\in D}\|f-Tf\|_{L^p(\mu,E)}<\eps .
\]
\end{theorem}

\begin{proof}
Fix \(\eps>0\).  Choose positive numbers \(\eta,\delta\) such that
\[
  \delta+(1+C)\eta<\eps .
\]
By \(L^p\)-weak-compact tightness, there is a weakly compact set \(W\subset E\) with
\[
  \sup_{f\in D}\left(\int_\Omega \dist(f(\omega),W)^p\,d\mu(\omega)\right)^{1/p}<\eta .
\]
Since \(E\) has the W.A.P. with constant \(C\), there is \(S\in W(E)\) such that
\[
  \|S\|\le C,\qquad \sup_{w\in W}\|w-Sw\|<\delta .
\]
Define \(T:L^p(\mu,E)\to L^p(\mu,E)\) by
\[
  (Tf)(\omega)=S(f(\omega)).
\]
Then \(T\) is bounded and \(\|T\|\le \|S\|\le C\).

We first estimate the approximation error.  For any \(x\in E\) and any \(w\in W\),
\[
  \|x-Sx\|
  \le \|x-w\|+\|w-Sw\|+\|S(w-x)\|
  \le \delta+(1+C)\|x-w\|.
\]
Taking the infimum over \(w\in W\) gives
\[
  \|x-Sx\|\le \delta+(1+C)\dist(x,W).
\]
Therefore, for \(f\in D\),
\[
\begin{aligned}
  \|f-Tf\|_{L^p}
  &\le \delta\,\mu(\Omega)^{1/p}
      +(1+C)\left(\int_\Omega \dist(f(\omega),W)^p\,d\mu(\omega)\right)^{1/p} \\
  &< \delta+(1+C)\eta<\eps .
\end{aligned}
\]

It remains to check that \(T\) is weakly compact.  Since \(S\) is weakly compact, the Davis-Figiel-Johnson-Pelczynski factorization theorem gives a reflexive Banach space \(R\) and operators \(A:E\to R\), \(B:R\to E\) such that \(S=BA\).  The induced Bochner operators
\[
  \widetilde A:L^p(\mu,E)\to L^p(\mu,R),\qquad
  \widetilde B:L^p(\mu,R)\to L^p(\mu,E)
\]
are bounded and satisfy \(T=\widetilde B\widetilde A\).  Because \(1<p<\infty\) and \(R\) is reflexive, \(L^p(\mu,R)\) is reflexive.  Hence \(T\) factors through a reflexive space and is weakly compact.
\end{proof}

\begin{corollary}\label{cor:reduction}
Let \(1<p<\infty\).  If \(E\) has the W.A.P. and every weakly compact subset of \(L^p([0,1];E)\) is \(L^p\)-weakly-compactly tight, then \(L^p([0,1];E)\) has the W.A.P.
\end{corollary}

\begin{proof}
Apply Theorem \ref{thm:tight-lift} to each weakly compact set \(D\subset L^p([0,1];E)\).
\end{proof}

\section{A concrete class: uniformly bounded \texorpdfstring{\(\delta\mathcal S\)}{delta-S} sets}

Rodriguez defines \(\delta\mathcal S\)-sets in Lebesgue-Bochner spaces: a set \(K\subset L^p(\mu,E)\) is a \(\delta\mathcal S\)-set if it is bounded and, for every \(\delta>0\), there is a weakly compact \(W\subset E\) such that
\[
  \mu(f^{-1}(W))\ge 1-\delta\qquad(f\in K).
\]
Every \(\delta\mathcal S\)-set is relatively weakly compact, but the converse fails in general.

\begin{corollary}\label{cor:bounded-deltas}
Let \(1<p<\infty\), and suppose \(E\) has the W.A.P.  Let \(D\subset L^p(\mu,E)\) be a weakly compact \(\delta\mathcal S\)-set which is uniformly essentially bounded: there is \(M<\infty\) such that
\[
  \|f(\omega)\|\le M\quad\text{for a.e. }\omega,\ \text{for every }f\in D.
\]
Then \(D\) is uniformly approximable by weakly compact operators on \(L^p(\mu,E)\), with the same W.A.P. constant as \(E\).
\end{corollary}

\begin{proof}
It suffices to show that \(D\) is \(L^p\)-weakly-compactly tight.  Fix \(\eta>0\).  By the \(\delta\mathcal S\) condition, choose a weakly compact \(W_0\subset E\) such that
\[
  \mu(f^{-1}(W_0))\ge 1-(\eta/M)^p
\]
for every \(f\in D\), with the evident interpretation if \(M=0\).  Put \(W=W_0\cup\{0\}\), still weakly compact.  Since \(\dist(f(\omega),W)=0\) on \(f^{-1}(W_0)\), and \(\dist(f(\omega),W)\le \|f(\omega)\|\le M\) almost everywhere,
\[
  \int_\Omega \dist(f(\omega),W)^p\,d\mu(\omega)
  \le M^p\,\mu(\Omega\setminus f^{-1}(W_0))<\eta^p
\]
uniformly over \(f\in D\).  Theorem \ref{thm:tight-lift} applies.
\end{proof}

\section{Limitations and novelty check}

This packet does not solve Problem 5.7 in full.  The missing step is not cosmetic: the literature on \(\delta\mathcal S\)-sets records that weak compactness in \(L^1(\mu,E)\), and likewise the natural \(L^p\)-range-control analogues, need not reduce to containment near functions whose values lie in a fixed weakly compact subset of \(E\).  Rodriguez explicitly notes that a strong \(L^p\)-neighborhood condition fails for \(X=L^1[0,1]\) and \(X=\ell^1\) when \(L^1(\mu)\) is infinite-dimensional.

The result should therefore be read as a lifting lemma and a reduction: any proof of \(L^p\)-weak-compact tightness for the weakly compact sets relevant to a given \(E\) immediately gives the W.A.P. for those sets in \(L^p(\mu,E)\).  The inner W.A.P. half of Problem 5.7 is not addressed here.

Bounded search performed before promotion: the run indexes were searched for arXiv:0309405, the title, ``Bochner'', ``weakly compact approximation'', ``inner W.A.P.'', and Problem 5.7.  Web searches for exact phrases including ``weakly compact approximation property Lp(X)'', ``weakly compact approximation Bochner'', and ``Problem 5.7 weakly compact approximation'' did not reveal a full later answer.  The source arXiv paper and Rodriguez's arXiv:1611.07199 were inspected locally.

\section{References}

\begin{itemize}
\item E. Odell and H.-O. Tylli, \emph{Weakly compact approximation in Banach spaces}, arXiv:math/0309405.
\item J. Rodriguez, \emph{A class of weakly compact sets in Lebesgue-Bochner spaces}, arXiv:1611.07199.
\item W. J. Davis, T. Figiel, W. B. Johnson, and A. Pelczynski, \emph{Factoring weakly compact operators}, J. Funct. Anal. 17 (1974), 311--327.
\end{itemize}

\end{document}
